% Generated draft: requires mathematical and source review.
\documentclass[11pt]{article}
\usepackage[margin=25mm]{geometry}
\usepackage{fontspec}
\setmainfont{Latin Modern Roman}
\usepackage{amsmath,amssymb,mathtools,unicode-math}
\setmathfont{Latin Modern Math}
\usepackage{enumitem,needspace,xurl,hyperref}
\setlength{\parindent}{0pt}
\setlength{\parskip}{5pt}
\setlength{\emergencystretch}{3em}
\newcommand{\N}{\mathbb{N}}
\newcommand{\Z}{\mathbb{Z}}
\newcommand{\Q}{\mathbb{Q}}
\newcommand{\R}{\mathbb{R}}
\newcommand{\C}{\mathbb{C}}
\newcommand{\F}{\mathbb{F}}
\newcommand{\A}{\mathbb{A}}
\newcommand{\PP}{\mathbb{P}}
\newcommand{\E}{\mathbb{E}}
\newcommand{\eps}{\varepsilon}
\newcommand{\dd}{\,\mathrm{d}}
\begin{document}
\section*{UnsolvedMath ID: 30001262}
\textbf{Source problem code:} catalogue entry
\subsection{Definitions and mathematical statement}
A group is \emph{elementary amenable} if it belongs to the smallest class containing all finite groups and all abelian groups and closed under taking subgroups, quotient groups, group extensions, and directed unions. A group \(G\) is \emph{virtually cyclic} if it contains a cyclic subgroup of finite index; finite groups are included in this definition.

Under the ordinary integral group-homological convention, \(G\) has type \(FP_1\) if the trivial left \(\Z G\)-module \(\Z\) admits a projective resolution
\[
\cdots \longrightarrow P_2\longrightarrow P_1\longrightarrow P_0\longrightarrow \Z\longrightarrow 0
\]
in which \(P_0\) and \(P_1\) are finitely generated. Ordinary type \(FP_\infty\) requires every \(P_i\) to be finitely generated.

The supplied question, interpreted using this ordinary integral convention, asks whether
\[
\forall G,\qquad
\bigl(G\text{ is elementary amenable and has type }FP_1\bigr)
\ \stackrel{?}{\Longrightarrow}\ 
\bigl(\exists C\leq G:\ C\text{ is cyclic and }[G:C]<\infty\bigr).
\]
Here \(G\) ranges over groups, with no further hypotheses supplied.

\textbf{Unresolved qualification.} The record writes only \(FP_1\), while its references and literature assessment raise a possible distinction between ordinary and Bredon homological finiteness. Bredon finiteness requires a specified family of subgroups and an associated orbit category; the record specifies neither. Consequently, the displayed ordinary interpretation is conditional and cannot be identified confidently with the intended Question 8. No Bredon family, coefficient convention for such a formulation, or additional group hypothesis is inferred. The supplied assertion concerning an \(FP_\infty\) theorem has the same unresolved finiteness-convention issue.
\subsection{Short English statement}
The supplied question asks whether every elementary amenable group satisfying the intended \(FP_1\) finiteness condition contains a cyclic subgroup of finite index. The meaning of \(FP_1\) requires review: ordinary integral finiteness is defined above conditionally, but the supplied Bredon reference may indicate a different intended condition. The supplied open-status label and claimed \(FP_\infty\) result are dated curation, not independently verified conclusions.
\subsection{Sources}
\begin{itemize}
\item Local problem registry: supplied complete record for permanent UnsolvedMath ID \(30001262\), source problem code \texttt{catalogue entry}, in the Oberwolfach Reports Open Problems collection. The record gives proposed year 2009 and names no individual proposer.
\item Supplied report citation: \emph{Manifold Perspectives} (2009), \emph{Oberwolfach Reports}, pp.\ 1487--1546, DOI \texttt{10.4171/owr/2009/27}. Source URL and literature-source label ``original OWR report'': \url{https://doi.org/10.4171/owr/2009/27}.
\item Supplied literature-source label ``report theorem and Question 8'': \url{https://ems.press/content/serial-article-files/46228}.
\item Supplied literature-source label ``Bredon finiteness paper'': \url{https://eprints.soton.ac.uk/349586/}. The record supplies no author, title, or theorem statement for this reference.
\item Literature assessment dated exactly \texttt{2026-08-21}: the supplied curation labels the problem open and states that the report proves the \(FP_\infty\) case, whereas Question 8 asks the \(FP_1\) case and is not settled by that theorem. This is supplied dated curation derived from the OWR material, not independently verified here; the status and cited claims require independent verification, including the intended finiteness conventions.
\item The registry labels the statement \texttt{corrected\_verified}; its supplied verification note describes retention of the mathematical wording with normalization of notation and punctuation. This label is preserved as source metadata and does not establish mathematical correctness of the formulation.
\end{itemize}
\end{document}
